\begin{theorem}[sigmaR-natCast]
\uses{sigmaR}
\label{sigmaR_natCast}
\lean{ArithmeticFunction.sigmaR_natCast}
% at /home/runner/work/trough/trough/library/PrimeNumberTheoremAnd/IwaniecKowalskiCh1.lean:363.0-379.29
\leanok
For natural exponents, $\sigma^R$ agrees with $\sigma$.
\end{theorem}
\begin{proof}
\leanok
The function $\sigma^R$ is defined as the sum of the $s$-th powers of the divisors of $n$. When
  $s$ is a natural number $k$, this definition coincides with the classical divisor power sum
  function $\sigma_k(n)$, which also sums the $k$-th powers of the divisors of $n$. Therefore, for
  natural exponents, we have $\sigma^R_k(n) = \sigma_k(n)$ when we view $\sigma_k(n)$ as a complex
  number. This can be shown by directly comparing the definitions and noting that both functions sum
  over the same set of divisors with the same exponentiation.
\end{proof}
